%%% Hlavní soubor. Zde se definují základní parametry a odkazuje se na ostatní části. %%%

%% Verze pro jednostranný tisk:
% Okraje: levý 40mm, pravý 25mm, horní a dolní 25mm
% (ale pozor, LaTeX si sám přidává 1in)
\documentclass[12pt,a4paper]{report}
\setlength\textwidth{145mm}
\setlength\textheight{247mm}
\setlength\oddsidemargin{15mm}
\setlength\evensidemargin{15mm}
\setlength\topmargin{0mm}
\setlength\headsep{0mm}
\setlength\headheight{0mm}
% \openright zařídí, aby následující text začínal na pravé straně knihy
\let\openright=\clearpage

%% Pokud tiskneme oboustranně:
% \documentclass[12pt,a4paper,twoside,openright]{report}
% \setlength\textwidth{145mm}
% \setlength\textheight{247mm}
% \setlength\oddsidemargin{15mm}
% \setlength\evensidemargin{0mm}
% \setlength\topmargin{0mm}
% \setlength\headsep{0mm}
% \setlength\headheight{0mm}
% \let\openright=\cleardoublepage

%% Použité kódování znaků: obvykle latin2, cp1250 nebo utf8:
\usepackage[utf8]{inputenc}

%% Ostatní balíčky
\usepackage{graphicx}
\usepackage{amsthm}
\usepackage{amssymb}
\usepackage{subcaption}

%% Balíček hyperref, kterým jdou vyrábět klikací odkazy v PDF,
%% ale hlavně ho používáme k uložení metadat do PDF (včetně obsahu).
%% POZOR, nezapomeňte vyplnit jméno práce a autora.
\usepackage[unicode]{hyperref}   % Musí být za všemi ostatními balíčky
\hypersetup{pdftitle=Parallelization of Clustering Algorithms}
\hypersetup{pdfauthor=Bc. Jakub Vlček}

%%% Drobné úpravy stylu

% Tato makra přesvědčují mírně ošklivým trikem LaTeX, aby hlavičky kapitol
% sázel příčetněji a nevynechával nad nimi spoustu místa. Směle ignorujte.
\makeatletter
\def\@makechapterhead#1{
  {\parindent \z@ \raggedright \normalfont
   \Huge\bfseries \thechapter. #1
   \par\nobreak
   \vskip 20\p@
}}
\def\@makeschapterhead#1{
  {\parindent \z@ \raggedright \normalfont
   \Huge\bfseries #1
   \par\nobreak
   \vskip 20\p@
}}
\makeatother

% Toto makro definuje kapitolu, která není očíslovaná, ale je uvedena v obsahu.
\def\chapwithtoc#1{
\chapter*{#1}
\addcontentsline{toc}{chapter}{#1}
}

\begin{document}

% Trochu volnější nastavení dělení slov, než je default.
\lefthyphenmin=2
\righthyphenmin=2

%%% Titulní strana práce

\pagestyle{empty}
\begin{center}

\large

Charles University in Prague

\medskip

Faculty of Mathematics and Physics

\vfill

{\bf\Large MASTER THESIS}

\vfill

\centerline{\mbox{\includegraphics[width=60mm]{img/logo.eps}}}

\vfill
\vspace{5mm}

{\LARGE Bc. Jakub Vlček}

\vspace{15mm}

% Název práce přesně podle zadání
{\LARGE\bfseries Parallelization of Clustering Algorithms}

\vfill

% Název katedry nebo ústavu, kde byla práce oficiálně zadána
% (dle Organizační struktury MFF UK)
Department of Software Engineering

\vfill

\begin{tabular}{rl}

Supervisor of the master thesis: & RNDr. Martin Kruliš, Ph.D. \\
\noalign{\vspace{2mm}}
Study programme: & Informatics \\
\noalign{\vspace{2mm}}
Specialization: & Software Systems \\
\end{tabular}

\vfill

% Zde doplňte rok
Prague 2015

\end{center}

\newpage

%%% Následuje vevázaný list -- kopie podepsaného "Zadání diplomové práce".
%%% Toto zadání NENÍ součástí elektronické verze práce, nescanovat.

%%% Na tomto místě mohou být napsána případná poděkování (vedoucímu práce,
%%% konzultantovi, tomu, kdo zapůjčil software, literaturu apod.)

\openright

\noindent
Dedication.

\newpage

%%% Strana s čestným prohlášením k diplomové práci

\vglue 0pt plus 1fill

\noindent
I declare that I carried out this master thesis independently, and only with the cited
sources, literature and other professional sources.

\medskip\noindent
I understand that my work relates to the rights and obligations under the Act No.
121/2000 Coll., the Copyright Act, as amended, in particular the fact that the Charles
University in Prague has the right to conclude a license agreement on the use of this
work as a school work pursuant to Section 60 paragraph 1 of the Copyright Act.

\vspace{10mm}

\hbox{\hbox to 0.5\hsize{%
In ........ date ............
\hss}\hbox to 0.5\hsize{%
Jakub Vlček
\hss}}

\vspace{20mm}
\newpage

%%% Povinná informační strana diplomové práce

\vbox to 0.5\vsize{
\setlength\parindent{0mm}
\setlength\parskip{5mm}

Název práce:
Parallelization of Clustering Algorithms
% přesně dle zadání

Autor:
Bc. Jakub Vlček

Katedra:  % Případně Ústav:
Katedra softwarového inženýrství 
% dle Organizační struktury MFF UK

Vedoucí diplomové práce:
RNDr. Martin Kruliš, Ph.D., Katedra softwarového inženýrství 
% dle Organizační struktury MFF UK, případně plný název pracoviště mimo MFF UK

Abstrakt:
V diplomové práci se zabývám využitím moderních procesorových architektur, zvláště pak grafických čipů, k paralelizaci náročných výpočetních problémů, jakým je například shluková analýza. Na tomto problému zkoumám možnosti zrychlení na různých typech architektur (CPU a GPU) a především závislosti různých vstupních dat a různých přístupů k paralelizaci. Ty se odvíjí například od počtu vlastností (dimenze), které zkoumané objekty mají. Dále je také v práci dbáno na co nejefektivnější využití daných architektur (propustnost paměti, využití co nejvyššího počtu jader, minimalizace závislostí), kde není problém pouze rozdíl mezi CPU a GPU, ale i v jednotlivých verzích konkrétní architektury.
% abstrakt v rozsahu 80-200 slov; nejedná se však o opis zadání diplomové práce

Klíčová slova:
dolování dat, shluková analýza, paralelizace, GPU, CUDA
% 3 až 5 klíčových slov

\vss}\nobreak\vbox to 0.49\vsize{
\setlength\parindent{0mm}
\setlength\parskip{5mm}

Title: Parallelization of Clustering Algorithms
% přesný překlad názvu práce v angličtině

Author:
Bc. Jakub Vlček

Department:
Department of Software Engineering
% dle Organizační struktury MFF UK v angličtině

Supervisor:
RNDr. Martin Kruliš, Ph.D., Department of Software Engineering
% dle Organizační struktury MFF UK, případně plný název pracoviště
% mimo MFF UK v angličtině

Abstract:
The thesis deals with the use of the newest processor architectures, especially graphics processors (GPU), for parallelization of sophisticated computational problems such as cluster analysis. On this problem, I investigate possibilities of speed-up different types of architectures (CPU and GPU) and especially the dependence of different input data and different approaches to parallelization. For example, problem is the diversity of input data (the number of object properties). Furthermore, this thesis deals with on the most efficient use of the architecture like memory bandwidth, use the maximum number of cores, minimizing dependencies. Problem is not only in difference between CPU and GPU architecture but also in the different versions of a particular architecture.
% abstrakt v rozsahu 80-200 slov v angličtině; nejedná se však o překlad
% zadání diplomové práce

Keywords:
data mining, cluster analysis, parallelization, GPU, CUDA
% 3 až 5 klíčových slov v angličtině

\vss}

\newpage

%%% Strana s automaticky generovaným obsahem diplomové práce. U matematických
%%% prací je přípustné, aby seznam tabulek a zkratek, existují-li, byl umístěn
%%% na začátku práce, místo na jejím konci.

\openright
\pagestyle{plain}
\setcounter{page}{1}
\tableofcontents

%%% Jednotlivé kapitoly práce jsou pro přehlednost uloženy v samostatných souborech
\pagestyle{plain}
\setcounter{page}{1}

\chapter{Úvod}
Hlavní motivací pro výběr této práce byl vysoký nárůst v poptávce na strojové zpracování dat, kde je velký důraz kladen na rychlost a efektivitu výpočtu. Často se zde setkáváme s protichůdnými požadavky na krátkou dobu výpočtu a~zpracování velkého množství dat.
Pokud pro takovéto úlohy využijeme moderní počítačové čipy podporující paralelní výpočty, můžeme výpočet urychlit rozdělením úlohy na menší nezávislé části, které pak lze řešit najednou. Hlavním problémem je právě toto rozdělení dat, protože neexistuje obecný způsob jeho provedení. Různé algoritmy totiž vyžadují jiné rozdělení dat a mnohdy závisí i~na vlastnostech dat samotných. Tuto problematiku je tedy nutné prozkoumat detailněji.\\

V této práci se zaměříme na možnosti zlepšení shlukové analýzy pomocí efektivní implementace, kde budeme klást důraz na maximální využití paralelních architektur, jako jsou například vícejádrové počítačové (CPU) a grafické (GPU) čipy. Budeme se zabývat především různými přístupy k paralelizaci algoritmu a závislosti na vlastnostech vstupních dat.
Data se od sebe velice liší, protože pochází z mnoha různých zdrojů, jako jsou 	například sociální sítě, finanční burzy, lékařské pozorování, průzkum vesmíru a mnoho dalších. Všechna však obsahují velmi užitečná data pro vědce, sociology, lékaře, ale i pro analytiky, makléře a například i pro cílenou reklamu.
Pokud se nám podaří urychlit shlukovou analýzu paralelizací výpočtu, budeme moci řešit problémy mnohem rychleji, což je pro některé odvětví zásadní. Naše řešení může dokonce dovolit využití shlukové analýzy i do odvětví, ve kterých byla dříve nepoužitelná ať už kvůli rychlosti výpočtu a nebo velkému objemu dat.\\

Protože shluková analýza zahrnuje spoustu možných přístupů k jejímu vy\-ře\-še\-ní, musíme si vybrat jeden konkrétní algoritmus, který budeme následně analyzovat. Pro tuto práci jsme zvolil algoritmus k-means a to především pro jeho široké využití, nízkou složitost a dobře známé matematické vlastnosti. Také jde o jeden z nejefektivnějších shlukových algoritmů zmiňovaných v literatuře~\cite{Aggarwal13}.\\

Algoritmus k-means můžeme definovat následovně: Vstupem je množina bodů a počet shluků $k$. Úkolem algoritmu je seskupit body do nepřekrývajících se shluků. Každý takový shluk navíc obsahuje speciální bod zvaný centroid, který je zkonstruován jako průměr ze všech bodů, kterým je mezi centroidy nejblíže. K výpočtu vzdálenosti lze použít datům odpovídající metriku. Nejčastěji se však využívá Eukleidovská vzdálenost~\cite{Zechner09}. Optimální množinu centroidů lze nalézt tak, že se snažíme minimalizovat součet vzdáleností mezi jednotlivými body a~středem daného shluku, který je reprezentován právě nejbližším centroidem. Je dokázáno, že vyřešit tento problém je NP-těžké už pro 2 shluky~\cite{Drineas04}. Algoritmus k-means optimální řešení pouze aproximuje postupnou konvergencí k lokálnímu minimu, které závisí na volbě počátečních centroidů.\\

Prvním krokem algoritmu k-means je výběr $k$ bodů, které poslouží jako počáteční centroidy. Velmi často se k tomuto účelu volí náhodné body ze vstupní množiny. V dalším kroku je pak každý bod přiřazen k nejbližšímu centroidu. Poté se pro nově vzniklé shluky spočítá nový centroid jako průměr obsažených bodů. Algoritmus pracuje iterativně opakováním těchto dvou kroků, dokud se centroidy neustálí (nezmění se mezi dvěma iteracemi) a nebo pokud není dosaženo předem určeného počtu iterací.\\

Protože je algoritmus tvořen mnoha datově paralelními úlohami, je dobré použít výpočetní prostředky, které jsou schopny této vlastnosti využít.
Dobrým kandidátem pro paralelizaci k-means jsou grafické karty~\cite{Zechner09} jakožto levné, masivně paralelní procesory. Obsahují tisíce jednoduchých jader, které jsou schopné počítat základní matematické operace. I přesto, že jejich hlavním zaměřením jsou primárně grafické úlohy, jsou tyto karty schopné poskytnout svůj výpočetní výkon i pro obecné úlohy (General-purpose computing on graphics processing units - GPGPU). Pokud se nám tedy podaří využít potenciálu grafických karet a zároveň datového paralelismu v algoritmu k-means k rovnoměrnému rozdělení výpočtů mezi jednotlivá výpočetní jádra GPU, budeme schopni výpočet mnohonásobně zrychlit.\\


Paralelizace algoritmu k-means ale přináší i úskalí, která je potřeba vyřešit. Algoritmus musí být optimalizován pro masivně paralelní prostředí. To především znamená, že musíme být schopni jednotlivé úkoly rovnoměrně rozdělit mezi tisíce výpočetních jader a postarat se o jejich efektivní využití. Pokud by se nám nepodařilo zapojit všechna dostupná jádra, ztratili bychom velké množství výkonu, což je samozřejmě nežádoucí. Další výzvou je pak co nejefektivnější využití paměťového modelu grafických karet, protože špatné používání paměti, jako na\-pří\-klad časté přístupy k nejpomalejším pamětím, mohou výkon algoritmu rapidně snížit.\\

V roce 2009 již zkoumal využití grafických karet pro algoritmus k-means Hong-Tao Bai a jeho týmem~\cite{Hong09}, ale tato práce byla spíše konceptem a neprozkoumala detailní možnosti paralelizace.
V naší práci se budeme soustředit na vlastnosti vstupních dat a jejich vliv na výkonnost algoritmu při odlišných přístupech k paralelizaci. Tento vztah je pro k-means rozhodující, protože právě vstupní data velice ovlivňují výpočet algoritmu a tím pádem i jeho efektivitu.
Dalším důvodem je pak velice rychlý vývoj technologií, tedy i grafických karet. Ty se od roku 2009 v mnohém posunuly kupředu, takže bychom se chtěli zaměřit i na výkon jednotlivých generací karet.\\

\hyperref[sec:clusteranalysis]{Kapitola~\ref*{sec:clusteranalysis}} obsahuje úvod do shlukové analýzy, různé druhy modelů shluků a shlukových algoritmů. \hyperref[sec:gpgpu]{Kapitola~\ref*{sec:gpgpu}} se zabývá GPGPU, CUDA frameworkem a~jeho analýzou. \hyperref[sec:implementation]{Kapitola~\ref*{sec:implementation}} obsahuje popis a analýzu naší implementace a paralelizačních metod. V \hyperref[sec:results]{kapitole~\ref*{sec:results}} je pak rozbor dosažených výsledků této práce a~efektivity paralelizace k-means.
\chapter{Implementation}
\section{Platform}
We choose a C++ language for implementing our solutions, because we need low level code optimization and features like SIMD instructions, which are not available in other languages. Even if the CUDA is developed to work with multiple languages, cooperation with  C++ is one of the best, because also CUDA kernels are written in C++. Also platform independence of C++ is useful because we developed programs on Windows and than run on Linux clusters.

\section{K-means algorithm} \label{sec:kMeansAlgorithm}
We focused on standard k-means algorithm~\cite{hartigan1979algorithm}, which input is set of $n$ points in $d$-dimensional space $R^d$ and an number of centroids $k$. The task is to assign points to centroids and minimize the distance between points and centroids. Algorithm works iteratively and each iteration consist of two steps. At first points are assigned to the nearest centroids and than new centroids are computed. Algorithm ends after known number of steps or when nothing changes between two iterations. Output of the this algorithm are centroids coordinates or points with assigned clusters.\\
As we can see in serial implementation~\autoref{lst:kMeansCode}, there are many cycles dependent on data size, so with growing number of data, serial version comes much slower. For example in two nested cycles on lines \autoref{lst:nestedCyclesBegin}-\autoref{lst:nestedCyclesEnd} of \autoref{lst:kMeansCode}, slowdown is great and this part of code is good target for parallelization.\\

\begin{lstlisting}[caption={k-means algorithm}, language=C++, escapeinside={\%*}{*)}, label={lst:kMeansCode}, numbers=left]
CreateLSH(width, hashTablesCount, basicHashFunctions)
{
  // Main iteration loop
  for( i = 0; i < iterationsCount; i++)
  {
    assign_to_clusters(data, means);
    compute_means(data_means);
  }
}

// Assign points to clusters
%*\label{lst:assignToClusters}*)assign_to_clusters(data, means)
{
//Iterate through data
%*\label{lst:nestedCyclesBegin}*)  for( d = 0; d < data_size; d++)
  {
    assigned_cluster = 0;
    min_distance = MAX_VALUE;
// Iterate through means
    for( m = 0; m < means_size; m++)
    {
      distance = count_distance( data[d], means[m] );
// If m-cluster is closer, reassign cluster
      if ( distance < min_distance )
      {
        assigned_cluster = m;
        min_distance = distance;
      }
    }
    data[d]->assigned_cluster = assigned_cluster;
%*\label{lst:nestedCyclesEnd}*)  }
}

// Compute new means
%*\label{lst:computeNewMeans}*)compute_means(data, means)
{
// Create new means array, all values are set to 0
  for( m = 0; m < means_size; m++)
  {
    new_means.push_back(new mean(0));
  }
  
// Iterate through data and accumulate data in appropriate new mean
  for( d = 0; d < data_size; d++)
  {
    assigned_cluster = data[d]->assigned_cluster;
    new_means[assigned_cluster]->contained_points++;
    for( dim = 0; dim < dimensions; dim++)
    {
      new_means[assigned_cluster]->coordinates[dim] +=
        data[d]->coordinates[dim];
    }
  }
  
// Count new means coordinates (arithmetic mean)
  for( m = 0; m < means_size; m++)
  {
    for( dim = 0; dim < dimensions; dim++)
    {
      new_means[m]->coordinates[dim] /=
        new_means[m]->contained_points;
    }
  }
}

// Function for counting distance
count_distance(point, mean)
{
  sum_squares = 0;
  for( dim = 0; dim < dimensions; dim++)
  {
    difference = point->coordinates[dim] - mean->coordinates[dim];
    sum_squares += difference * difference;
  }
  return sqrt( sum_squares );
}
\end{lstlisting}

\section{Parallelization}
As we mentioned in \autoref{sec:kMeansAlgorithm}, classical k-means algorithm contains many compute-intensive cycles, which could be parallelized. Problem is that the algorithm contains some compute dependencies, for example parallelization of nested cycles on lines \autoref{lst:nestedCyclesBegin}-\autoref{lst:nestedCyclesEnd} of \autoref{lst:kMeansCode} is problematic, because we need to accumulate coordinates in shared variable. These dependencies must be solved some way for correct parallelization.

\subsection{CPU versions}
We tried to used the most modern technologies offered by the newest CPU architectures to speed up k-means algorithm by parallelization. We created two levels of parallelization, first based on thread level (using more logical cores on single CPU) and than vectorization level using SIMD instructions.

\subsubsection{Thread level parallelization}
For thread level parallelization, we choose Threading Building Blocks (TBB) library from Intel. This library offers algorithms and data structures for parallel programing using multi-core processors but relieve programmer from complications arising from the use of native threading packages and handling particular threads. Manual thread creation, execution, synchronization and termination is hidden by using by creating TBB tasks which are processed on individual cores dynamically by the library. Also efficient use of CPU caches is automated.\\
We reworked program so points are split in groups, each group forming one TBB task.\\
\begin{description}
\item[Assigning to clusters]The first part of algorithm, assigning to clusters~\autoref{lst:kMeansCode},~\autoref{lst:assignToClusters}, is not so problematic, because only write operation is assigning cluster to point. Each point is contained in one and only one TBB task, which is processed by single thread, so it could not be read or written from other threads.\\
Second possible solution is to split means into groups and parallelize algorithm by tasks containing means. This solution is significantly slower, because we must synchronize writing information to each point about assigned cluster. If we want to remember this information in means (by accumulating assigned points to each mean) we hit a problem, because after assigning point to mean, we could not know if other mean is not closer to processed point so this way o parallelization is too computationally intensive.\\
The second part of the algorithm, counting new means~\autoref{lst:kMeansCode},~\autoref{lst:computeNewMeans}, is more problematic, because we approach means from different TBB tasks, so accumulating of coordinates must be synchronized. We solved this problem by local copy of means in each TBB task, so we avoided the collisions. For joining tasks, we used TBB reduce algorithm, when to tasks are joined, local arrays of means coordinate sums are correctly added.\\
Also in this part of algorithm exists second possible solution by approaching parallelization by splitting means. This means that for each mean, we must iterate through all points and accumulate coordinates, which means 
\end{description}
\subsubsection{Vectorization}
\subsection{GPU versions}
\section{Tools}
\chapter{Measurements and results} \label{sec:results}

\section{Measurements}

\section{Results}


% Ukázka použití některých konstrukcí LateXu (odkomentujte, chcete-li)
% %%% Ukázka použití některých konstrukcí LaTeXud

\subsection{Ukázka \LaTeX{}u}
\label{ssec:ukazka}

This short subsection serves as an~example of basic \LaTeX{} constructs,
which can be useful for writing a~thesis.

Let us start with lists:

\begin{itemize}
\item The logo of Matfyz is displayed in figure~\ref{fig:mff}.
\item This is subsection~\ref{ssec:ukazka}.
\item Citing literature~\cite{lamport94}.
\end{itemize}

Different kinds of dashes:
red-black (short),
pages 16--22 (middle),
$45-44$ (minus),
and this is --- as you could have expected --- a~sentence-level dash,
which is the longest.
(Note that we have follwed \verb|a| by a~tilde instead of a~space
to avoid line breaks at that place.)

\newtheorem{theorem}{Theorem}
\newtheorem*{define}{Definition}	% Definice nečíslujeme, proto "*"

\begin{define}
A~{\sl Tree} is a connected graph with no cycles.
\end{define}

\begin{theorem}
This theorem is false.
\end{theorem}

\begin{proof}
False theorems do not have proofs.
\end{proof}

\begin{figure}
	\centering
	\includegraphics[width=30mm]{../img/logo.eps}
	\caption{Logo of MFF UK}
	\label{fig:mff}
\end{figure}


\chapter*{Conclusion}
\addcontentsline{toc}{chapter}{Conclusion}


%%% Seznam použité literatury
%%% Seznam použité literatury je zpracován podle platných standardů. Povinnou citační
%%% normou pro diplomovou práci je ISO 690. Jména časopisů lze uvádět zkráceně, ale jen
%%% v kodifikované podobě. Všechny použité zdroje a prameny musí být řádně citovány.

\def\bibname{Bibliography}
\begin{thebibliography}{99}
\addcontentsline{toc}{chapter}{\bibname}

%\bibitem{lamport94}
%  {\sc Lamport,} Leslie.
%  \emph{\LaTeX: A Document Preparation System}.
%  2. vydání.
%  Massachusetts: Addison Wesley, 1994.
%  ISBN 0-201-52983-1.

\bibitem{EstivillCastro02}
{\sc Estivill-Castro,} Vladimir
\emph{Why so many clustering algorithms.}
[online]. p.~65-75. DOI 10.1145/568574.568575.

\bibitem{Zechner09}
{\sc Zechner,} Mario, {\sc Granitzer, } Michael
\emph{Accelerating K-Means on the Graphics Processor via CUDA}
First International Conference on Intensive Applications and Services
2013
ISBN: 978-1466558212 

\bibitem{Aggarwal13}
{\sc Aggarwal,} Charu C., {\sc Reddy, } Chandan K.
\emph{Data Clustering: Algorithms and Applications}
Chapman \& Hall/CRC Data Mining and Knowledge Discovery Series
2009
ISBN: 978-1-4244-3683-5

\bibitem{Drineas04}
{\sc Drineas,} P., {\sc Frieze,} A., {\sc Kannan,} R., {\sc Vempala,}  S., {\sc Vinay,}  V.
\emph{Clustering large graphs via the singular value decomposition}
Machine Learning., 56(1-3):9–33.
2004
ISSN: 0885-6125

\bibitem{Bottou95}
{\sc Bottou,} L\'{e}on, {\sc Bengio,} Yoshua
\emph{Convergence Properties of the K-Means Algorithms}
Advances in Neural Information Processing Systems 7
p.~585-592
MIT Press, 1995

\bibitem{Kim14}
{\sc Kim,} Kyoungok
\emph{Voronoi Cell-Based Clustering Using a Kernel Support.}
Knowledge and Data Engineering, IEEE Transactions
ISSN: 1041-4347
DOI 10.1109/TKDE.2014.2359662. 

\bibitem{Vallender73}
{\sc Vallender,} S.S.
\emph{Calculation of the Wasserstein Distance Between Probability Distributions on the Line.}
[online]. s.~784--786. DOI: 10.1137/1118101.

\bibitem{Beecks10}
{\sc Beecks,} Christian, {\sc Uysal,} Merih Seran, {\sc Seidl,} Thomas
\emph{Signature Quadratic Form Distance.}
Proceedings of the ACM International Conference on Image and Video Retrieval - CIVR , 2010.
ACM
ISBN 978-1-4503-0117-6
DOI 10.1145/1816041.1816105

\bibitem{Tan05}
{\sc Tan,} Pang-Ning, {\sc Michael,} Steinbach, {\sc Kumar,} Vipin
\emph{Introduction to data mining.}
Pearson New International Edition
Boston: Pearson Addison Wesley, 2005
ISBN 978-0-3213-2136-7

\bibitem{Sun15}
{\sc Sun,} Xili, {\sc Tian,} Shoucai, {\sc Lu,} Yonggang
\emph{High dimensional data clustering by partitioning the hypergraphs using dense subgraph partition}
MIPPR 2015: Pattern Recognition and Computer Vision
DOI 10.1117/12.2205743.

\bibitem{Kirk12}
{\sc Kirk,}David B., {\sc Hwu}Wen-mei W.:
\emph{Programming Massively Parallel Processors}
Second Edition: A Hands-on Approach
2012
ISBN: 0124159923 

\bibitem{Sanders10}
{\sc Sanders,} Jason, {\sc Kandrot,} Edward
\emph{CUDA by Example: An Introduction to General-Purpose GPU Programming}
NVIDIA
2010
ISBN: 0-13-138768-5 

\bibitem{Scarpino11}
{\sc Scarpino,} Matthew
\emph{OpenCL in Action: How to Accelerate Graphics and Computations}
Manning Publications
2011
ISBN: 1617290173 

\bibitem{Hong09}
{\sc Bai,} Hong-Tao, {\sc He,} Li-li, {\sc Ouyang,} Dan-tong, {\sc Li,} Zhan-shan, {\sc Li, } He
\emph{K-means on commodity GPUs with CUDA}
 Computer Science and Information Engineering
 2009
 WRI World Congress on. Volume 3.
 IEEE (2009) 651–655 

\bibitem{CUDAGuide}
Nvidia
\emph{CUDA C PROGRAMMING GUIDE}
Nvidia, 2015
available online: http://docs.nvidia.com/cuda/cuda-c-programming-guide/
\end{thebibliography}

%%% Tabulky v diplomové práci, existují-li.
\chapwithtoc{List of Tables}

%%% Použité zkratky v diplomové práci, existují-li, včetně jejich vysvětlení.
\chapwithtoc{List of Abbreviations}

%%% Přílohy k diplomové práci, existují-li (různé dodatky jako výpisy programů,
%%% diagramy apod.). Každá příloha musí být alespoň jednou odkazována z vlastního
%%% textu práce. Přílohy se číslují.
\chapwithtoc{Attachments}

\openright
\end{document}
